\documentclass[10pt,a4paper]{amsart}
\usepackage{amsmath,amssymb,amsthm,graphicx}
% Retained minimal preamble for the Stacks Project chapter varieties.tex.
% Verbatim excerpt of the chapter preamble: Xy-pic setup (native preamble lines 14-18), the native theorem
% declarations with the mathematical macros (native preamble lines 154-217), and the project own no-op
% definitions of the reference/slogan/history display environments (native wrapper.tex lines 1-3) that the
% counted statement uses for its slogan. Copied mechanically; nothing added, removed or redefined.
% For commutative diagrams we use Xy-pic
\usepackage[all]{xy}

% We use 2cell for 2-commutative diagrams.
\xyoption{2cell}
% Theorem environments.
%
\theoremstyle{plain}
\newtheorem{theorem}[subsection]{Theorem}
\newtheorem{proposition}[subsection]{Proposition}
\newtheorem{lemma}[subsection]{Lemma}

\theoremstyle{definition}
\newtheorem{definition}[subsection]{Definition}
\newtheorem{example}[subsection]{Example}
\newtheorem{exercise}[subsection]{Exercise}
\newtheorem{situation}[subsection]{Situation}

\theoremstyle{remark}
\newtheorem{remark}[subsection]{Remark}
\newtheorem{remarks}[subsection]{Remarks}

\numberwithin{equation}{subsection}

% Macros
%
\def\lim{\mathop{\mathrm{lim}}\nolimits}
\def\colim{\mathop{\mathrm{colim}}\nolimits}
\def\Spec{\mathop{\mathrm{Spec}}}
\def\Hom{\mathop{\mathrm{Hom}}\nolimits}
\def\Ext{\mathop{\mathrm{Ext}}\nolimits}
\def\SheafHom{\mathop{\mathcal{H}\!\mathit{om}}\nolimits}
\def\SheafExt{\mathop{\mathcal{E}\!\mathit{xt}}\nolimits}
\def\Sch{\mathit{Sch}}
\def\Mor{\mathop{\mathrm{Mor}}\nolimits}
\def\Ob{\mathop{\mathrm{Ob}}\nolimits}
\def\Sh{\mathop{\mathit{Sh}}\nolimits}
\def\NL{\mathop{N\!L}\nolimits}
\def\CH{\mathop{\mathrm{CH}}\nolimits}
\def\proetale{{pro\text{-}\acute{e}tale}}
\def\etale{{\acute{e}tale}}
\def\QCoh{\mathit{QCoh}}
\def\Ker{\mathop{\mathrm{Ker}}}
\def\Im{\mathop{\mathrm{Im}}}
\def\Coker{\mathop{\mathrm{Coker}}}
\def\Coim{\mathop{\mathrm{Coim}}}

% Boxtimes
%
\DeclareMathSymbol{\boxtimes}{\mathbin}{AMSa}{"02}

%
% Macros for moduli stacks/spaces
%
\def\QCohstack{\mathcal{QC}\!\mathit{oh}}
\def\Cohstack{\mathcal{C}\!\mathit{oh}}
\def\Spacesstack{\mathcal{S}\!\mathit{paces}}
\def\Quotfunctor{\mathrm{Quot}}
\def\Hilbfunctor{\mathrm{Hilb}}
\def\Curvesstack{\mathcal{C}\!\mathit{urves}}
\def\Polarizedstack{\mathcal{P}\!\mathit{olarized}}
\def\Complexesstack{\mathcal{C}\!\mathit{omplexes}}
% \Pic is the operator that assigns to X its picard group, usage \Pic(X)
% \Picardstack_{X/B} denotes the Picard stack of X over B
% \Picardfunctor_{X/B} denotes the Picard functor of X over B
\def\Pic{\mathop{\mathrm{Pic}}\nolimits}
\def\Picardstack{\mathcal{P}\!\mathit{ic}}
\def\Picardfunctor{\mathrm{Pic}}
\def\Deformationcategory{\mathcal{D}\!\mathit{ef}}
\newenvironment{reference}{}{}
\newenvironment{slogan}{}{}
\newenvironment{history}{}{}
\usepackage{cleveref}
\crefname{originallemma}{original lemma}{original lemmas}
\Crefname{originallemma}{Original lemma}{Original lemmas}
\makeatletter
\newlabel{lemma-geometrically-normal}{{038O}{source}{}{sourcepointer-1}{}}
\newlabel{lemma-geometrically-normal@cref}{{[originallemma][0][]038O}{[1][0][]source}{}{}{}}
\makeatother
\crefname{originallemma}{original lemma}{original lemmas}
\Crefname{originallemma}{Original lemma}{Original lemmas}
\makeatletter
\newlabel{lemma-normalization-locally-algebraic}{{0BXR}{source}{}{sourcepointer-2}{}}
\newlabel{lemma-normalization-locally-algebraic@cref}{{[originallemma][0][]0BXR}{[1][0][]source}{}{}{}}
\makeatother
\crefname{originallemma}{original lemma}{original lemmas}
\Crefname{originallemma}{Original lemma}{Original lemmas}
\makeatletter
\newlabel{limits-lemma-descend-closed-immersion-finite-presentation}{{01ZP}{source}{}{sourcepointer-3}{}}
\newlabel{limits-lemma-descend-closed-immersion-finite-presentation@cref}{{[originallemma][0][]01ZP}{[1][0][]source}{}{}{}}
\makeatother
\crefname{originallemma}{original lemma}{original lemmas}
\Crefname{originallemma}{Original lemma}{Original lemmas}
\makeatletter
\newlabel{limits-lemma-descend-finite-finite-presentation}{{01ZO}{source}{}{sourcepointer-4}{}}
\newlabel{limits-lemma-descend-finite-finite-presentation@cref}{{[originallemma][0][]01ZO}{[1][0][]source}{}{}{}}
\makeatother
\crefname{originallemma}{original lemma}{original lemmas}
\Crefname{originallemma}{Original lemma}{Original lemmas}
\makeatletter
\newlabel{limits-lemma-descend-finite-presentation}{{01ZM}{source}{}{sourcepointer-5}{}}
\newlabel{limits-lemma-descend-finite-presentation@cref}{{[originallemma][0][]01ZM}{[1][0][]source}{}{}{}}
\makeatother
\crefname{originallemma}{original lemma}{original lemmas}
\Crefname{originallemma}{Original lemma}{Original lemmas}
\makeatletter
\newlabel{limits-lemma-descend-opens}{{01Z4}{source}{}{sourcepointer-6}{}}
\newlabel{limits-lemma-descend-opens@cref}{{[originallemma][0][]01Z4}{[1][0][]source}{}{}{}}
\makeatother
\crefname{originallemma}{original lemma}{original lemmas}
\Crefname{originallemma}{Original lemma}{Original lemmas}
\makeatletter
\newlabel{morphisms-lemma-birational-isomorphism-over-dense-open}{{0BAJ}{source}{}{sourcepointer-7}{}}
\newlabel{morphisms-lemma-birational-isomorphism-over-dense-open@cref}{{[originallemma][0][]0BAJ}{[1][0][]source}{}{}{}}
\makeatother
\title{Geometric normality of the normalization is reached after a finite purely inseparable field extension}
\author{The Stacks Project Authors}
\date{stacks-project:varieties.tex:local-snapshot-20261011; GNU FDL 1.2 or later}
\begin{document}
\maketitle
\noindent\textit{Editorial note.} All delivered mathematical text is the author's own, copied from the pinned Stacks Project chapter file varieties.tex; no statement or proof step was rewritten, reordered or removed.\par\medskip
\noindent\textit{Editorial note.} The last sentence of the proof is the author's own 'It follows readily from this that nu is the normalization morphism'; the sentence is reproduced verbatim, and the two properties it refers to were both established inside the counted span.\par\medskip
\noindent\textit{Editorial note.} The counted spans carry 7 navigation stubs and no bibliography citation.\par\medskip
\noindent\textit{Editorial note.} The preamble is the chapter's retained minimal preamble (sources/stacks-varieties/retained-preamble.tex lines 1-77), a verbatim excerpt of the native preamble.\par\medskip
\noindent\textit{Editorial note.} The label lemma-finite-extension-geometrically-normal is defined only in varieties.tex, so no cross-chapter restatement had to be weighed for this card.\par\medskip
\noindent\textit{Allowed standard prerequisites (editorial).}\begin{itemize}
\item Normalization of a scheme locally of finite type over a field. Same chapter, Lemma lemma-normalization-locally-algebraic, invoked at proof lines 5001-5002. Supplies the normalization morphism over the perfect closure.
\item Descent of finite presentation, finiteness, open subschemes and closed immersions along a filtered colimit. Limits, Lemmas limits-lemma-descend-finite-presentation, limits-lemma-descend-finite-finite-presentation, limits-lemma-descend-opens and limits-lemma-descend-closed-immersion-finite-presentation, invoked at proof lines 5003-5005, 5009-5010 and 5018-5021. Supplies each of the objects and properties descended to a finite subextension.
\item Geometric normality is preserved when passing to a larger base field. Same chapter, Lemma lemma-geometrically-normal, invoked at proof lines 5007-5008. Supplies that the descended morphism is geometrically normal over k'.
\item A birational morphism is an isomorphism over a dense open. Morphisms, Lemma morphisms-lemma-birational-isomorphism-over-dense-open, invoked at proof lines 5015-5017. Supplies the dense open on which the normalization is a closed immersion.
\end{itemize}\medskip
\ifdefined\hypertarget\hypertarget{sourcepointer-1}{}\fi
\noindent\textit{Source reference.} Original lemma 038O: sources/stacks-varieties/pinned-original-source.tex, source lines 2160--2160. Pointer outside the counted span, invoked once for context; the load-bearing argument of the counted statement is the author's proof reproduced in full in the proof span.\par\medskip
\ifdefined\hypertarget\hypertarget{sourcepointer-2}{}\fi
\noindent\textit{Source reference.} Original lemma 0BXR: sources/stacks-varieties/pinned-original-source.tex, source lines 4903--4903. Pointer outside the counted span, invoked once for context; the load-bearing argument of the counted statement is the author's proof reproduced in full in the proof span.\par\medskip
\ifdefined\hypertarget\hypertarget{sourcepointer-3}{}\fi
\noindent\textit{Source reference.} Original lemma 01ZP: sources/stacks-chow/tags.txt, source lines 2602--2602. Pointer outside the counted span, invoked once for context; the load-bearing argument of the counted statement is the author's proof reproduced in full in the proof span.\par\medskip
\ifdefined\hypertarget\hypertarget{sourcepointer-4}{}\fi
\noindent\textit{Source reference.} Original lemma 01ZO: sources/stacks-chow/tags.txt, source lines 2601--2601. Pointer outside the counted span, invoked once for context; the load-bearing argument of the counted statement is the author's proof reproduced in full in the proof span.\par\medskip
\ifdefined\hypertarget\hypertarget{sourcepointer-5}{}\fi
\noindent\textit{Source reference.} Original lemma 01ZM: sources/stacks-chow/tags.txt, source lines 2599--2599. Pointer outside the counted span, invoked once for context; the load-bearing argument of the counted statement is the author's proof reproduced in full in the proof span.\par\medskip
\ifdefined\hypertarget\hypertarget{sourcepointer-6}{}\fi
\noindent\textit{Source reference.} Original lemma 01Z4: sources/stacks-chow/tags.txt, source lines 2579--2579. Pointer outside the counted span, invoked once for context; the load-bearing argument of the counted statement is the author's proof reproduced in full in the proof span.\par\medskip
\ifdefined\hypertarget\hypertarget{sourcepointer-7}{}\fi
\noindent\textit{Source reference.} Original lemma 0BAJ: sources/stacks-chow/tags.txt, source lines 14212--14212. Pointer outside the counted span, invoked once for context; the load-bearing argument of the counted statement is the author's proof reproduced in full in the proof span.\par\medskip
\begin{lemma}
\label{lemma-finite-extension-geometrically-normal}
Let $k$ be a field. Let $X$ be an algebraic $k$-scheme.
Then there exists a finite purely inseparable extension $k'/k$
such that the normalization $Y$ of $X_{k'}$ is geometrically normal over $k'$.
\end{lemma}

\begin{proof}
Let $K = k^{perf}$ be the perfect closure. Let $Y_K$ be the
normalization of $X_K$, see Lemma \ref{lemma-normalization-locally-algebraic}.
By Limits, Lemma \ref{limits-lemma-descend-finite-presentation}
there exists a finite sub extension $K/k'/k$ and a morphism
$\nu : Y \to X_{k'}$ of finite presentation whose base change to $K$
is the normalization morphism $\nu_K : Y_K \to X_K$.
Observe that $Y$ is geometrically normal over $k'$
(Lemma \ref{lemma-geometrically-normal}).
After increasing $k'$ we may assume $Y \to X_{k'}$ is finite
(Limits, Lemma \ref{limits-lemma-descend-finite-finite-presentation}).
Since $\nu_K : Y_K \to X_K$ is the normalization morphism,
it induces a birational morphism $Y_K \to (X_K)_{red}$.
Hence there is a dense open $V_K \subset X_K$ such that
$\nu_K^{-1}(V_K) \to V_K$ is a closed immersion
(inducing an isomorphism of $\nu_K^{-1}(V_K)$ with $V_{K, red}$, see
for example Morphisms, Lemma
\ref{morphisms-lemma-birational-isomorphism-over-dense-open}).
After increasing $k'$ we find $V_K$ is the base change of a dense open
$V \subset Y$ and the morphism $\nu^{-1}(V) \to V$ is a closed immersion
(Limits, Lemmas \ref{limits-lemma-descend-opens} and
\ref{limits-lemma-descend-closed-immersion-finite-presentation}).
It follows readily from this that $\nu$ is the normalization
morphism and the proof is complete.
\end{proof}

\end{document}
